

\medskip
\noindent\SGAInlineRefTarget{sga3:iii:remark:4-8-0}\textbf{Remark 4.8.0.}\footnote{Editor's note: This remark,
used in the proposition that follows, has been inserted here; it appeared
in \SGARefLink{sga3:iii:paragraph:4-10}{4.10} of the original.}
\textit{Let \(S\) be a scheme, \(S'\) a closed subscheme of \(S\), \(X\) an
\(S\)-scheme, \(Y\) an \(S\)-subscheme of \(X\), and put
\(X'=X\times_S S'\) and \(Y'=Y\times_S S'\). Then there is a
surjective morphism of \(\mathscr O_{Y'}\)-modules}
\[
 \mathscr N_{Y/X}\otimes_{\mathscr O_Y}\mathscr O_{Y'}
 \twoheadrightarrow \mathscr N_{Y'/X'}.
\]

Indeed, after replacing \(X\) by a suitable open subscheme, one may
suppose that \(Y\) is closed and is defined by an ideal
\(\mathscr I_Y\) of \(\mathscr O_X\). The image of \(\mathscr I_Y\) in
\(\mathscr O_{X'}\) is then the ideal \(\mathscr I_{Y'}\) defining
\(Y'\), and there is a surjective morphism of
\(\mathscr O_{Y'}\)-modules
\[
 \pi:
 (\mathscr I_Y/\mathscr I_Y^2)
 \otimes_{\mathscr O_Y}\mathscr O_{Y'}
 \twoheadrightarrow
 \mathscr I_{Y'}/\mathscr I_{Y'}^2.
\]
Suppose in addition that
\(\mathscr O_Y=\mathscr O_X/\mathscr I_Y\) is flat over
\(\mathscr O_S\). The natural morphism
\[
 \mathscr I_Y\otimes_{\mathscr O_X}\mathscr O_{X'}
 \longrightarrow \mathscr I_{Y'}
\]
is then bijective (cf. EGA IV\textsubscript{2}, 2.1.8).

There is therefore the following commutative diagram with exact rows:

\begin{center}
 \SGAThreeNativeInput[0.96\textwidth]
   {native_diagrams/exp3/diagram_039_exp3_p052_remark480_exact_rows.tex}
\end{center}

It follows from the snake lemma that
\begin{equation}
 \mathscr N_{Y/X}\otimes_{\mathscr O_Y}\mathscr O_{Y'}
 \xrightarrow{\sim}\mathscr N_{Y'/X'}
 \qquad\text{if \(Y\) is flat over \(S\).}
 \tag{4.8.0}\label{sga3:iii:remark:4-8-0:equation:4-8-0}
\end{equation}

\medskip
\noindent\SGAInlineRefTarget{sga3:iii:proposition:4-8}\textbf{Proposition 4.8.}
\textit{Let \(S,S_0,S_{\mathscr J}\), and
\(\mathscr I,\mathscr J\) be as in \SGARefLink{sga3:iii:proposition:4-5}{4.5}.}
\textit{Let \(X\) be an \(S\)-scheme, let \(Y\) be a subscheme of
\(X\), and let \(i:Y\hookrightarrow X\) be the immersion.}
\begin{enumerate}
\item[\SGAInlineRefTarget{sga3:iii:proposition:4-8:item:i}(i)] \textit{For every \(S\)-morphism \(f:T\to X\) such that
  \(f_{\mathscr J}:T_{\mathscr J}\to X_{\mathscr J}\) factors through
  \(Y_{\mathscr J}\), one can define an obstruction}
  \begin{equation}
   c(X,Y,f)\in
   \operatorname{Hom}_{\mathscr O_{T_0}}
   \left(
    f_0^*\bigl(\mathscr N_{Y/X}
      \otimes_{\mathscr O_Y}\mathscr O_{Y_0}\bigr),
    \mathscr J\mathscr O_T
   \right)
   \tag{*}\label{sga3:iii:proposition:4-8:equation:star}
  \end{equation}
  \textit{whose vanishing is equivalent to the existence of a
  factorization of \(f\) through \(Y\).}

\item[\SGAInlineRefTarget{sga3:iii:proposition:4-8:item:ii}(ii)] \textit{Let \(Y'\) be a second subscheme of \(X\). Suppose
  that \(Y'_{\mathscr J}=Y_{\mathscr J}\) and that \(Y\) and \(Y'\) are
  flat over \(S\). Then there are isomorphisms (cf. \SGARefLink{sga3:iii:remark:4-8-0}{4.8.0})}
  \[
   \mathscr J\mathscr O_Y
   \simeq
   \mathscr J\otimes_{\mathscr O_{S_0}}\mathscr O_{Y_0}
   \simeq
   \mathscr J\mathscr O_{Y'},
   \qquad
   \mathscr N_{Y/X}\otimes_{\mathscr O_Y}
     \mathscr O_{Y_{\mathscr J}}
   \xrightarrow{\sim}
   \mathscr N_{Y_{\mathscr J}/X_{\mathscr J}},
  \]
  \textit{and hence an isomorphism}
  \[
   \begin{split}
   u:\ &\operatorname{Hom}_{\mathscr O_{Y_0}}
   \left(
    \mathscr N_{Y_{\mathscr J}/X_{\mathscr J}}
      \otimes_{\mathscr O_{Y_{\mathscr J}}}\mathscr O_{Y_0},
    \mathscr J\mathscr O_Y
   \right)\\
   &\xrightarrow{\sim}
   \operatorname{Hom}_{\mathscr O_{Y_0}}
   \left(
    \mathscr N_{Y/X}\otimes_{\mathscr O_Y}\mathscr O_{Y_0},
    \mathscr J\mathscr O_{Y'}
   \right).
   \end{split}
  \]
  \textit{Writing \(i':Y'\hookrightarrow X\) for the canonical
  immersion and \(d(Y,Y')\) for the deviation of \SGARefLink{sga3:iii:remark:4-5-1}{4.5.1}, one has}
  \begin{equation}
   c(X,Y,i')=u\bigl(d(Y,Y')\bigr).
   \tag{**}\label{sga3:iii:proposition:4-8:equation:double-star}
  \end{equation}

\item[\SGAInlineRefTarget{sga3:iii:proposition:4-8:item:iii}(iii)] \textit{The canonical morphism}
  \[
   \mathscr N_{Y/X}\xrightarrow{D}i^*(\Omega^1_{X/S})
  \]
  \textit{(cf. SGA 1, II, formula 4.3)}
  \textit{induces a morphism}
  \[
   D_0:
   \mathscr N_{Y/X}\otimes_{\mathscr O_Y}\mathscr O_{Y_0}
   \longrightarrow
   \Omega^1_{X_0/S_0}\otimes_{\mathscr O_{X_0}}\mathscr O_{Y_0},
  \]
  \textit{and hence, for every \(S\)-morphism \(f:T\to X\) as in
  \SGARefLink{sga3:iii:proposition:4-8:item:i}{(i)}, a morphism}
  \[
   \begin{split}
   v_{f_0}:\ &
   \operatorname{Hom}_{\mathscr O_{T_0}}
   \left(f_0^*(\Omega^1_{X_0/S_0}),\mathscr J\mathscr O_T\right)
   \\[-0.15em]
   &\longrightarrow
   \operatorname{Hom}_{\mathscr O_{T_0}}
   \left(
    f_0^*\bigl(\mathscr N_{Y/X}
      \otimes_{\mathscr O_Y}\mathscr O_{Y_0}\bigr),
    \mathscr J\mathscr O_T
   \right),\\
   &\hspace{8em}a\longmapsto a\circ f_0^*(D_0),
   \end{split}
  \]
  \textit{where the first group above is
  \(\operatorname{Hom}_{X^+}(T,L_X)\), cf. \SGARefLink{sga3:iii:definition:0-1-5}{0.1.5}. For
  \(a\in\operatorname{Hom}_{X^+}(T,L_X)\), one has}
  \begin{equation}
   c(X,Y,a\cdot f)-c(X,Y,f)=v_{f_0}(a),
   \tag{***}\label{sga3:iii:proposition:4-8:equation:double-starstar}
  \end{equation}
  \textit{where \(a\cdot f\) denotes the composite morphism}
  \[
   T\xrightarrow{a\times f}L_X\times_{X^+}X\longrightarrow X.
  \]
\end{enumerate}

We shall prove part (i) of the proposition, leaving it to the reader to
(not) verify assertions (ii) and (iii); this verification proceeds by
reduction to the affine case and then by comparison of the explicit
definitions.

Let us therefore prove \SGARefLink{sga3:iii:proposition:4-8:item:i}{(i)}. The morphism \(f:T\to X\) defines a
morphism of sheaves of rings\footnote{Editor's note: On the one hand,
the hypothesis that \(\mathscr I\) be nilpotent---that is, that \(X_0\)
have the same underlying topological space as \(X\)---has been deleted;
on the other hand, the sentence that follows has been expanded.}
\[
 \phi:\mathscr O_X\longrightarrow f_*(\mathscr O_T).
\]
Let \(U\) be an open subscheme of \(X\) in which \(Y\) is closed. Since
\(T\), respectively \(Y_{\mathscr J}\), has the same underlying space
as \(T_{\mathscr J}\), respectively \(Y\), the continuous map underlying
\(f\) sends \(T\) into \(U\). Since \(U\) is open in \(X\), the morphism
\(\phi\) induces a morphism of sheaves of rings
\[
 \mathscr O_U=\mathscr O_X|_U\longrightarrow f_*(\mathscr O_T),
\]
that is, \(f\) factors through \(U\).

We may therefore restrict to the case in which \(Y\) is closed, hence
defined by a sheaf of ideals \(\mathscr I_Y\). The morphism \(f\)
factors through \(Y\) if and only if the composite
\[
 \mathscr I_Y\longrightarrow\mathscr O_X
 \longrightarrow f_*(\mathscr O_T)
\]
is zero. Since \(f_{\mathscr J}\) factors through
\(Y_{\mathscr J}\), the composite
\[
 \mathscr I_{Y_{\mathscr J}}\longrightarrow
 \mathscr O_{X_{\mathscr J}}\longrightarrow
 f_*(\mathscr O_{T_{\mathscr J}})
\]
is zero. Consider the following commutative diagram, whose first row is
exact:

\begin{center}
 \SGAThreeNativeInput[0.82\textwidth]
   {native_diagrams/exp3/diagram_040_exp3_p053_prop48_obstruction_diagram.tex}
\end{center}

It follows that \(\phi\) maps \(\mathscr I_Y\) into
\(f_*(\mathscr J\mathscr O_T)\).\footnote{Editor's note: The argument
that follows has been expanded.} Since \(\mathscr J^2=0\), the
\(\mathscr O_X\)-module \(f_*(\mathscr J\mathscr O_T)\), with its
\(\mathscr O_X\)-module structure induced by \(\phi\), is annihilated by
\(\mathscr I_Y\). Consequently, \(\phi\) induces a morphism of
\(\mathscr O_X\)-modules
\[
 h:i_*(\mathscr N_{Y/X})
   =\mathscr I_Y/\mathscr I_Y^2
 \longrightarrow f_*(\mathscr J\mathscr O_T).
\]

On the other hand, one has the following cartesian squares:

\begin{center}
 \SGAThreeNativeInput[0.46\textwidth]
   {native_diagrams/exp3/diagram_041_exp3_p054_prop48_cartesian_squares.tex}
\end{center}

Here \(\tau_{T_0},\tau_{X_0},\tau_{Y_0}\) are the closed immersions
obtained by base change from \(S_0\hookrightarrow S\). Since
\(\mathscr J\mathscr O_T\) is a quasi-coherent
\(\mathscr O_T\)-module annihilated by \(\mathscr I\), there is an
isomorphism
\[
 \mathscr J\mathscr O_T
 \simeq
 (\tau_{T_0})_*\tau_{T_0}^*(\mathscr J\mathscr O_T),
\]
and hence
\[
 f_*(\mathscr J\mathscr O_T)
 \simeq
 (\tau_{X_0})_*(f_0)_*\tau_{T_0}^*(\mathscr J\mathscr O_T).
\]
Thus, by adjunction, \(h\) corresponds to a morphism of
\(\mathscr O_{T_0}\)-modules
\[
 h_0:
 f_0^*\tau_{X_0}^*i_*(\mathscr N_{Y/X})
 \longrightarrow
 \tau_{T_0}^*(\mathscr J\mathscr O_T).
\]
\smallskip
\noindent{\footnotesize\textit{Translator's note.} The French prose and this \(h_0\)-display print \(i_{T_0}\) for the
closed immersion labeled \(\tau_{T_0}\) in the cartesian square and the
preceding base-change formulas. We use \(\tau_{T_0}\) consistently.}

Now
\[
 \tau_{X_0}^*i_*(\mathscr N_{Y/X})
 \simeq
 (i_0)_*\tau_{Y_0}^*(\mathscr N_{Y/X})
 =\mathscr N_{Y/X}\otimes_{\mathscr O_Y}\mathscr O_{Y_0}.
\]
Returning to the customary abuse of notation
\(\tau_{T_0}^*(\mathscr J\mathscr O_T)=\mathscr J\mathscr O_T\), the
morphism \(h_0\) identifies with a morphism of
\(\mathscr O_{T_0}\)-modules
\[
 h_0:
 f_0^*\bigl(\mathscr N_{Y/X}
   \otimes_{\mathscr O_Y}\mathscr O_{Y_0}\bigr)
 \longrightarrow\mathscr J\mathscr O_T,
\]
which is the required obstruction \(c(X,Y,f)\). This proves \SGARefLink{sga3:iii:proposition:4-8:item:i}{(i)}.

When \(f\) is the immersion \(i':Y'\hookrightarrow X\), one sees that
\(c(X,Y,i')\) comes from the morphism
\[
 \mathscr I_Y\lhook\joinrel\longrightarrow\mathscr O_X
 \longrightarrow\mathscr O_{Y'},
\]
and hence corresponds, by \SGARefLink{sga3:iii:complement:4-4-1}{4.4.1} and \SGARefLink{sga3:iii:remark:4-5-1}{4.5.1}, to the class \(d(Y,Y')\).
This proves (ii).

\smallskip
\noindent{\footnotesize\textit{Translator's note.} At this point the French proof prints \(i':Y\hookrightarrow X\), while part \SGARefLink{sga3:iii:proposition:4-8:item:ii}{(ii)}
defines \(i'\) as the immersion of \(Y'\). The latter domain is used
here.}

Let us prove (iii). First,
\(D:\mathscr N_{Y/X}\to i^*(\Omega^1_{X/S})\) induces a morphism
\[
 \begin{split}
 D_0:\quad
 \tau_{Y_0}^*(\mathscr N_{Y/X})
 &\longrightarrow
 \tau_{Y_0}^*i^*(\Omega^1_{X/S})\\
 &=i_0^*\tau_{X_0}^*(\Omega^1_{X/S}).
 \end{split}
\]
Since \(X_0=X\times_S S_0\), one has
\[
 \tau_{X_0}^*(\Omega^1_{X/S})\simeq\Omega^1_{X_0/S_0}
 \qquad\text{(cf. EGA IV\textsubscript{4}, 16.4.5).}
\]
This gives the asserted morphism
\[
 D_0:
 \mathscr N_{Y/X}\otimes_{\mathscr O_Y}\mathscr O_{Y_0}
 \longrightarrow
 \Omega^1_{X_0/S_0}\otimes_{\mathscr O_{X_0}}\mathscr O_{Y_0}.
\]
We shall verify equality \((***)\) after the following remark.

\medskip
\noindent\SGAInlineRefTarget{sga3:iii:remark:4-9}\textbf{Remark 4.9.}
The obstruction \(c(X,Y,f)\) is computed locally on \(T\). Let
\(U=\operatorname{Spec}(C)\) be an affine open subset of \(T\) lying
over an affine open subset \(\operatorname{Spec}(A)\) of \(X\), itself
lying over an affine open subset \(\operatorname{Spec}(\Lambda)\) of
\(S\). Let \(J\subset I\subset\Lambda\), respectively
\(I_Y\subset A\), be the ideals corresponding to
\(\mathscr J\subset\mathscr I\), respectively \(\mathscr I_Y\), put
\(B=A/I_Y\), and let \(\phi:A\to C\) be the morphism of
\(\Lambda\)-algebras corresponding to \(f:T\to X\). Since
\(f(T_{\mathscr J})\subset Y_{\mathscr J}\), one has
\(\phi(I_Y)\subset JC\), and \(\phi\) therefore induces a morphism of
\(\Lambda\)-algebras
\[
 B\longrightarrow C/JC\longrightarrow C_0=C/IC.
\]
The obstruction \(c=c(X,Y,f)\) is then computed by the following
commutative diagram:

\begin{center}
 \SGAThreeNativeInput[0.60\textwidth]
   {native_diagrams/exp3/diagram_042_exp3_p054_remark49_local_obstruction.tex}
\end{center}

That is, over \(U\), it is the unique element of
\[
 \operatorname{Hom}_{C_0}
 \left(I_Y/I_Y^2\otimes_B C_0,JC\right)
 =
 \operatorname{Hom}_{B_0}
 \left(I_Y/I_Y^2\otimes_B B_0,JC\right)
\]
such that, with the evident notation,
\[
 c(\overline x\otimes_B1)=\phi(x)
 \qquad\text{for every \(x\in I_Y\).}
\]

\footnote{Editor's note: What follows has been added.}
We can now complete the proof of \SGARefLink{sga3:iii:proposition:4-8:item:iii}{4.8(iii)}. Equality \((***)\) may be
checked locally on \(T\), so we are reduced to the affine situation
described above. Let \(d_{A/\Lambda}:A\to\Omega^1_{A/\Lambda}\) denote
the differential. Over \(U\), the element \(a\) corresponds to an
element \(a_U\) of
\[
 \begin{split}
 &\operatorname{Hom}_{C_0}
 \left(\Omega^1_{A_0/\Lambda_0}\otimes_{A_0}C_0,JC\right)
 \\[-0.15em]
 &\qquad\simeq
 \operatorname{Hom}_{B_0}
 \left(\Omega^1_{A/\Lambda}\otimes_A B_0,JC\right)
 \simeq
 \operatorname{Hom}_{A}(\Omega^1_{A/\Lambda},JC).
 \end{split}
\]
On the one hand, \(v_{f_0}(a)\) corresponds over \(U\) to
\(a_U\circ D_0\), where \(D_0\) is the morphism of
\(B\)-modules
\[
 \begin{aligned}
 I_Y/I_Y^2&\longrightarrow\Omega^1_{A/\Lambda}\otimes_A B_0,\\
 x+I_Y^2&\longmapsto d_{A/\Lambda}(x)\otimes1.
 \end{aligned}
\]
On the other hand (cf. the proofs of \SGARefLink{sga3:iii:theorem:0-1-8}{0.1.8} and \SGARefLink{sga3:iii:corollary:0-1-9}{0.1.9}), the morphism of
\(\Lambda\)-algebras \(\phi':A\to C\) corresponding to \(a\cdot f\)
differs from \(\phi\) by the \(\Lambda\)-derivation \(A\to JC\)
associated with \(a_U\); that is,
\[
 \phi'=\phi+a_U\circ d_{A/\Lambda}
 =\phi+a_U\circ(d_{A/\Lambda}\otimes1).
\]
Consequently, writing \(c'=c(X,Y,a\cdot f)\), for every \(x\in I_Y\),
and writing \(\overline x\) for its image in \(I_Y/I_Y^2\), one has
\[
 \begin{split}
 (c'-c)(\overline x\otimes1)
 &=a_U\bigl(d_{A/\Lambda}(x)\otimes1\bigr)\\
 &=(a_U\circ D_0)(\overline x)
 =v_{f_0}(a)(\overline x).
 \end{split}
\]
This proves that \(c'-c=v_{f_0}(a)\).

\medskip
\noindent\SGAInlineRefTarget{sga3:iii:paragraph:4-10}\textbf{4.10.}
\SGARefLink{sga3:iii:paragraph:4-10}{Remark 4.10} of the original has been suppressed, since it was made
obsolete by the addition of \SGARefLink{sga3:iii:remark:4-8-0}{Remark 4.8.0}.

